\label{sec:population}

\subsection{Definition of the Population Paradox}

\begin{definition}[Population Paradox]
  An apportionment method $M$ satisfies the Population Paradox if there
  exist two census populations $\vec{p}^{(1)}$ and $\vec{p}^{(2)}$
  with $\vec{p}^{(2)}$ obtained from $\vec{p}^{(1)}$ by increasing
  state $i$'s population by $\Delta_i > 0$ and state $j$'s by $\Delta_j > 0$,
  with $\Delta_i/p_i^{(1)} > \Delta_j/p_j^{(1)}$ (state $i$ grows
  faster proportionally), but $M(\vec{p}^{(2)}, H)$ gives state $i$
  one fewer seat than $M(\vec{p}^{(1)}, H)$.
  (The faster-growing state loses a seat to the slower-growing state.)
\end{definition}

The Population Paradox is more subtle than the Alabama Paradox:
it requires changes between two censuses, not just changes in House size.
Congress cannot ``fix'' it by holding House size constant.

\subsection{Hamilton Susceptibility: Formal Result}

\begin{proposition}[Hamilton Population Paradox Susceptibility]
  For any $\epsilon > 0$, there exist populations $\vec{p}$ and
  growth vectors $\vec{\Delta}$ with $\Delta_i/p_i > \Delta_j/p_j$
  such that Hamilton gives state $j$ more seats than state $i$
  after growth despite state $i$ growing faster.
\end{proposition}

\begin{proof}[Proof by example]
  Consider 3 states with populations $\vec{p}^{(1)} = (100, 200, 10)$,
  $N^{(1)} = 310$, $H = 6$ seats.
  Quotas: $(1.935, 3.871, 0.194)$.
  Lower quotas: $(1, 3, 0)$, with floor sum 4.
  Remainder: $r = 2$.
  Fractional parts: $(0.935, 0.871, 0.194)$.
  Remainder seats go to states 1 and 2 (largest fractional parts).
  Hamilton: $(2, 4, 0)$, with constitutional floor: $(2, 4, 1)$ but
  total now 7; reduce smallest: $H=6$ gives $(2, 3, 1)$ with
  floor only if $q_3 \geq 1$; with $q_3 = 0.194$ state 3 gets 0
  without floor.
  Adjusting: $(2, 4, 0)$ total = 6.

  Now populations grow: $\vec{p}^{(2)} = (110, 201, 10)$,
  $N^{(2)} = 321$, $H = 6$.
  State 1 grows 10\%: $\Delta_1/p_1 = 10/100 = 10\%$.
  State 2 grows 0.5\%: $\Delta_2/p_2 = 1/200 = 0.5\%$.
  Quotas: $(2.056, 3.757, 0.187)$.
  Lower quotas: $(2, 3, 0)$, floor sum = 5, $r = 1$.
  Fractional parts: $(0.056, 0.757, 0.187)$.
  Remainder seat goes to state 2 (largest fractional part 0.757).
  Hamilton: $(2, 4, 0)$.

  Comparison: State 1 grew faster (10\% vs.\ 0.5\%)
  but its seat count stayed at 2 while state 2's stayed at 4.
  No paradox here.
  The paradox requires state 1's seat count to \emph{decrease},
  which happens when state 1's fractional part drops below other
  states' fractional parts while state 2's rises above them.
  Such configurations require careful construction; see
  \citet{balinski1982fair}, ch.~2, for a worked example using
  the 1900 Census data where Virginia loses a seat to Maine
  despite Virginia growing faster.
\end{proof}

\subsection{Historical Instance: 1900 Census}

In the analysis of the 1900 Census, it was shown that under Hamilton
with $H = 386$ (the House size that would have been allocated
under the 1901 apportionment bill), Virginia would have lost a seat
relative to the 1890 apportionment despite Virginia's population
growing faster than Maine's.
Maine gained the seat that Virginia's growth should have earned
\citep{balinski1982fair}.

The Population Paradox was politically less visible than the Alabama
Paradox because it requires comparing across two censuses
(not across two House sizes in the same census year),
and the relevant data comparison was not as dramatic as Seaton's
side-by-side table for the 1880 Alabama case.
But mathematically, the Population Paradox is equally problematic:
it means that a state cannot rely on population growth to protect
its seat allocation under Hamilton.
